% sec-05: Document 05, the San Diego Municipal Code update for the La Jolla site.
% STAGE 3.  Three column widths retuned; the permit table's Relationship column
% widened so "Parallel with 4" cannot wrap to two lines.
\docpart{San Diego Municipal Code}{San Diego Municipal Code Update - PDAC LLM
Oncology Clinical Trial Site Requirements, La Jolla Community Plan Area}

\section{Purpose and Applicability}

The purpose of this update is to establish the land use status, the permit
pathway, and the operating standards for a pancreatic oncology clinical trial
site incorporating an LLM-directed clinical workflow and robotic procedure
suites, within the La Jolla Community Plan area of the City of San Diego. It
applies to a facility that holds or has applied for a state authorization under
document 01 of this package and that does not hold a general acute care hospital
license.

The site is identified in this package by community plan area and by zone. The
street address is withheld until a lease is executed, because publishing a
parcel before terms are agreed prices the parcel. Nothing in this document
should be read as a representation that any specific parcel has been secured;
none has \cite{sdmc}.

\section{Zoning and Land Use}

\begin{mvtable}
\begin{tabularx}{\textwidth}{>{\raggedright\arraybackslash}p{2.7cm} >{\raggedright\arraybackslash}p{2.2cm} >{\raggedright\arraybackslash}p{2.8cm} >{\raggedright\arraybackslash}X}
\headrow \thead{Zone} & \thead{Use status} & \thead{Permit} & \thead{Condition attached by this update}\\
\midrule
Commercial office & Permitted & Neighborhood development permit & Robot envelopes wholly interior; no exterior mechanical addition above the parapet \\
Commercial community & Permitted & Neighborhood development permit & Loading and waste pickup limited to the hours in §5 \\
Mixed use, commercial ground floor & Conditionally permitted & Conditional use permit, process 3 & Acoustic separation from any residential use above or adjacent \\
Residential, any density & Not permitted & None available & Not permitted in any residential zone in the community plan area \\
Open space and coastal bluff & Not permitted & None available & Not permitted \\
\end{tabularx}
\end{mvtable}
\tabcapl{tab:sdzoning}{Zone by zone. Two zones permit the use outright, one
permits it conditionally, and two do not permit it at all. A use that is not
permitted anywhere is stated as such rather than omitted.}

The use is classified as a medical office and outpatient procedure use with an
accessory research function. It is not classified as a hospital, and the site
does not provide inpatient beds, an emergency department, or overnight stays. A
participant whose condition requires inpatient care is transferred under
document 11 §2.

\section{Conditional Use Permit Findings}

Where a conditional use permit is required, the hearing officer shall make each
of the following findings.

\begin{enumerate}
\item The proposed use is consistent with the La Jolla Community Plan and the
  General Plan.

\item The proposed use will not adversely affect the applicable land use plan.

\item The proposed use will not be detrimental to the public health, safety, and
  welfare.

\item The proposed use complies with the regulations of the Land Development
  Code, including any allowable deviations pursuant to the Code.

\item Every robotic operating envelope, as defined in document 09 §3, lies
  wholly within the building envelope, and no robot instance operates in a
  location visible from, or reachable from, the public right of way.

\item The applicant has filed a model governance record consisting of the model
  version register required by document 03 §3, and has agreed to make it
  available to the City on request in connection with any complaint.
\end{enumerate}

Findings (e) and (f) are specific to this use. The first four are the findings
the City already makes for any conditional use, and are restated here so that a
hearing officer reading this document alone has the complete set.

\section{Coastal Overlay and Community Plan Consistency}

Much of the La Jolla Community Plan area lies within the Coastal Overlay Zone.
The following applies where a proposed site does.

\begin{enumerate}
\item A tenant improvement wholly inside an existing permitted structure, with
  no increase in building height, no increase in gross floor area, no change to
  the exterior envelope, and no reduction in public access, is an improvement
  that does not intensify a coastal resource impact.

\item The applicant shall nonetheless file for a coastal development permit
  where the City determines that the change of use, as distinct from the
  construction, triggers one under the Coastal Act \cite{coastalact}. This
  document does not assert that no permit is required; it states the analysis
  the applicant must complete and file.

\item Review under the California Environmental Quality Act proceeds on the
  existing facilities exemption where the project involves no expansion of use
  beyond that existing at the time of the determination, and on a categorical
  exemption for minor alterations otherwise. The applicant shall state on the
  record which exemption is relied on and shall file the supporting analysis
  \cite{ceqa}.

\item Where a proposed site lies outside the Coastal Overlay Zone, subdivisions
  (a) through (c) do not apply and the ordinary permit sequence in §6 governs.
\end{enumerate}

\section{Operating Standards}

\begin{mvtable}
\begin{tabularx}{\textwidth}{>{\raggedright\arraybackslash}p{3.5cm} >{\raggedright\arraybackslash}p{4.6cm} >{\raggedright\arraybackslash}X}
\headrow \thead{Standard} & \thead{Limit} & \thead{Basis and cross-reference}\\
\midrule
Hours of operation & 06:00 to 20:00 weekdays; 08:00 to 14:00 Saturday; closed Sunday & Front matter site table; documents 12 §1 and 14 §4 use the same window \\
Sunday and after-hours exception & A procedure already in progress, or a participant safety event & Document 11 §1 \\
Noise at the property line & 55 dBA daytime, 50 dBA evening & Mechanical equipment is the only exterior source; the robots are interior \\
Exterior lighting & Full cutoff, 3000 K maximum, off by 22:00 except at exits & Coastal and residential adjacency \\
Delivery window & 07:00 to 17:00 weekdays & Loading geometry in document 10 §4 \\
Regulated medical waste pickup & Twice weekly, Tuesday and Friday, 07:00 to 09:00 & Waste streams in document 09 §5 \\
Exterior queuing & Prohibited; all waiting is interior & Document 10 §4 \\
\end{tabularx}
\end{mvtable}
\tabcapl{tab:sdoperating}{Seven operating standards. The hours are the single
most consequential row: this is an extended-day site, not the twenty-four-hour
facility described in the San Francisco predecessor, and the difference follows
from a protocol that schedules visits rather than accepting them on demand.}

\section{Permits, Fees, and Plan Check}

The permit sequence below determines the move-in schedule in document 14 §6,
which is why the parallel and sequential relationships are stated here rather
than assumed.

\begin{mvtable}
\begin{tabularx}{\textwidth}{>{\raggedright\arraybackslash}p{0.9cm} >{\raggedright\arraybackslash}p{4.4cm} >{\raggedright\arraybackslash}p{2.5cm} >{\raggedright\arraybackslash}X}
\headrow \thead{Step} & \thead{Permit or review} & \thead{Relationship} & \thead{Depends on}\\
\midrule
1 & Zoning verification & First & Nothing \\
2 & Neighborhood development or conditional use permit & Sequential & Step 1 \\
3 & Coastal development permit, if required & Parallel with 4 & Step 1 \\
4 & Building plan check & Sequential & Step 2 \\
5 & Mechanical, electrical, plumbing & Parallel with 4 & Step 2 \\
6 & Fire plan check & Parallel with 4 & Step 2 \\
7 & Hazardous materials business plan & Parallel with 4 & Step 2 \\
8 & County health plan check & Sequential & Steps 4 through 7 \\
9 & Certificate of occupancy & Last & Step 8 and commissioning under document 08 §6 \\
\end{tabularx}
\end{mvtable}
\tabcapl{tab:sdpermits}{Nine permits in sequence. Four of the nine run in
parallel, which is the only reason the forty-six-week move-in schedule in
document 14 §6 closes; a strictly sequential reading of the same nine adds
roughly eleven weeks.}
